\studytypesubsection{LLMs for Synthesis}
\label{sec:llms-for-synthesis}

In the synthesis role, LLMs integrate and interpret information from multiple sources to produce higher-level findings such as themes, patterns, or conceptual frameworks. They can also generate synthetic content (e.g., source code, bug-fix pairs, requirements) for downstream training or evaluation.

\studytypeparagraph{Description}

Unlike annotation (see \annotators), which focuses on categorizing or labeling individual data points, synthesis refers to the process of integrating and interpreting information from multiple sources to generate higher-level insights, identify patterns across datasets, and develop conceptual frameworks or theories. LLMs may be able to support synthesis tasks in SE research by processing and distilling information from qualitative data sources. 
Although synthesis in the preceding notion refers to abstraction and interpretation across multiple data sources, the term is sometimes also used to refer to generating synthetic content (e.g., source code, bug-fix pairs, or requirements) that is then used in downstream tasks to train, fine-tune, or evaluate existing models or tools.
In this case, the synthesis is done primarily using the LLM and its training data; the input is limited to basic instructions and examples.

\studytypeparagraph{Examples} 

Published examples of applying LLMs for synthesis in SE remain scarce. However, some recent work in other domains is instructive~\cite{DBLP:journals/ase/BanoHZT24}.
% GROUNDING DBLP:journals/ase/BanoHZT24: "the integration of Large Language Models (LLMs) like ChatGPT into qualitative research in software engineering, much like in other professional domains"
\citeauthor{DBLP:conf/wsese/BarrosANKKNB25} conducted a systematic mapping study on using LLMs for qualitative research~\cite{DBLP:conf/wsese/BarrosANKKNB25}. The included studies span fields such as healthcare, education, and cultural studies, where LLMs supported qualitative methods, such as grounded theory and thematic analysis, by aiding in pattern identification.
% GROUNDING DBLP:conf/wsese/BarrosANKKNB25: "The studies exploring the use of LLMs as a tool to support qualitative data analysis are concentrated in various fields, including healthcare, education, cultural studies, and analysis of technological applications."
In SE, \citeauthor{DBLP:conf/wsese/LecaVSS25} explored how LLMs have been applied for qualitative data analysis (QDA) and proposed general strategies and guidelines for their application~\cite{DBLP:conf/wsese/LecaVSS25}. \citeauthor{DBLP:journals/corr/abs-2511-14528} complemented this perspective by studying the opportunities and limitations of introducing LLM-based support into QDA, and by formulating recommendations for embedding human–AI collaboration across the thematic analysis phases~\cite{DBLP:journals/corr/abs-2511-14528}. Building on these, subsequent work has proposed hybrid frameworks combining LLM support with human-led QDA. \citeauthor{DBLP:journals/corr/abs-2402-01386} designed an LLM-driven multi-agent system that integrates AI with human decision-making to automate qualitative data analysis methods~\cite{DBLP:journals/corr/abs-2402-01386}. Their system generated initial codes, developed themes, and summarized text. Similarly, \citeauthor{DBLP:journals/corr/abs-2510-18456} compared the performance of humans and LLMs in coding, theme development, definition, and refinement, creating guidelines for a hybrid-LLM framework~\cite{DBLP:journals/corr/abs-2510-18456}.
% GROUNDING DBLP:conf/wsese/LecaVSS25: "This study aimed to investigate how LLMs are currently used in qualitative analysis and how they can be used in software engineering research"
% GROUNDING DBLP:journals/corr/abs-2511-14528: "A reflective workshop with 25 ISERN researchers guided participants through structured discussions of LLM-assisted open coding, theme generation, and theme reviewing, using color-coded canvases to document perceived opportunities, limitations, and recommendations."
% GROUNDING DBLP:journals/corr/abs-2402-01386: "we utilized LLMs to automate and expedite the qualitative data analysis processes. We employed a multi-agent model, where each agent was tasked with executing distinct, individual research related activities"
% GROUNDING DBLP:journals/corr/abs-2510-18456: "We systematically compared human and LLM-generated codes and themes between March and July 2025."
Finally, \citeauthor{DBLP:journals/corr/abs-2506-21138}'s work is an example of using LLMs to create synthetic datasets.
They presented an approach to generate synthetic requirements, showing that they \enq{can match or surpass human-authored requirements for specific classification tasks}~\cite{DBLP:journals/corr/abs-2506-21138}.
% GROUNDING DBLP:journals/corr/abs-2506-21138: "our results establish that synthetic requirements can match or surpass human-authored requirements for specific classification tasks"
